\scanpage{59}\mspage{58}
% slanted
\margin{alors $\mathcal F'$ satisfait aussi $C1$, $C2$, et $\unsure{M'|M}$ \deleted{\unsure{s'identifie}} \inserted{\unsure{désignant|désigné|figuré|défini} par} ses objets irréductibles, la relation d'ordre (de $M'\unsure{)|{}}$ induite par $\unsure{\overline{\mathcal F}'|\mathcal F'}$ étant celle \deleted{\unsure{d'}}induit par $M$. Si $\mathcal F$ satisfait $C3$, alors $\mathcal F'$ aussi.}\margin{\begin{tikzpicture}[x=1cm,y=1cm]\draw[<->] (.1,.55) .. controls (.7,.5) and (.1,.1) .. (.8,0);\draw (-.1,-.18)--(1.45,-.18);\end{tikzpicture}}

\deleted{\mbox{\uline{C 3}} Soient $F,G,H$ \deleted{trois} figures entre elles compatibles. \mbox{\inserted{i.e. deux à deux \unsure{en paires.|en paire.}}} Alors $\{F,G,H\}$ est une \unsure{configuration|configur.} i.e. \deleted{$F,G,H$} $\exists$ figure $K$ telle que $F,G,H$ sont des sous figures de $K$ (donc existe borne supérieure de}

\uline{C 3} Soient $F,G\in\mathcal F$, telles que
\[
\forall\,x\in\overline F,\ \overleftarrow{y}\in\overline G,\quad x\text{ et }y\text{ compatibles}
\]
\margin{\begin{tikzpicture}[x=1cm,y=1cm]\draw (0,0)--(0,2.15);\draw (.1,-.02)--(.1,2.12);\end{tikzpicture}}Alors $F$ et $G$ sont compatibles

% slanted
\margin{\uline{Cor} Pour que $F,G$ disjointes, il f et s. que $\forall x\in\overline F,y\in\overline G$, $x$ et $y$ le soient}

\uline{NB} L'inverse évident : Si $F$ et $G$ sont compatibles, alors pour toutes sous figures $F'$ de $F$ (p. ex. une ss-fg. él.) et $G'$ de $G$ (\unsure{id—|idem}) $F'$ et $G'$ le sont

\uline{Cor} Soient $F$ et $G,H$ \deleted{deux} figures, \deleted{\unsure{Si}} avec $F,G$ compatibles. Pour que $H$ soit : compatible : $F\cup G$, il f et s qu'il le soit : $F$ et : $G$. \deleted{\unsure{Deux|Il}}

Plus gén, soit $(F_i)_{i\in I}$ une configuration, $F$ la \deleted{\unsure{leur|borne}} \inserted{figure} réunion. Pour que $H$ soit compatible : $F$, il f et s que $H$ soit compatible \unsure{aux|avec} $F_i$
